\section{Results}
\label{sec:results}

This section reports fixed-policy results for Qwen2.5 and Qwen3 families
across two consumer-GPU setups: desktop RTX4090 and RTX5090-laptop.
We focus on paired speedup $S$, acceptance dynamics ($\alpha$, $B_{\text{eff}}$),
and token-timing metrics (time to first token, TTFT; time per output token,
TPOT), then summarize portability implications.

\begin{table*}[t]
\centering
\caption{Qwen2.5 fixed-policy results on RTX4090 (target: 3B, draft: 0.5B) against regime-matched autoregressive baselines.
Speedup is computed as a ratio of means against the corresponding baseline.}
\label{tab:headline-4090-qwen25}
\small
\begin{tabular}{l c c c c c c}
\hline
\textbf{Config} & $S$ & $\alpha$ & $B_{\text{eff}}$ & $T_{\text{mean}}$ (s) & TTFT (ms) & TPOT (ms) \\
\hline
0.5B, $k{=}4$, det   & \textbf{3.0674} & 0.4212 & 1.66 & 3.7210 & 102.35 & 29.06 \\
0.5B, $k{=}4$, stoch & 2.9394 & 0.4080 & 1.61 & 3.9631 & 105.26 & 30.89 \\
0.5B, $k{=}8$, det   & 2.4542 & 0.3515 & 2.72 & 4.6508 & 162.76 & 36.73 \\
0.5B, $k{=}8$, stoch & 2.4063 & 0.3444 & 2.67 & 4.8411 & 165.95 & 38.04 \\
0.5B, $k{=}16$, det  & 1.8050 & 0.2539 & 3.76 & 6.3234 & 277.43 & 49.34 \\
0.5B, $k{=}16$, stoch& 1.8185 & 0.2516 & 3.73 & 6.4059 & 278.85 & 49.81 \\
\hline
\end{tabular}
\end{table*}

\begin{table*}[t]
\centering
\caption{Qwen3 fixed-policy results on RTX4090 against regime-matched autoregressive baselines.
Speedup is computed as a ratio of means against the corresponding baseline.}
\label{tab:headline-4090}
\small
\begin{tabular}{l c c c c c c}
\hline
\textbf{Config} & $S$ & $\alpha$ & $B_{\text{eff}}$ & $T_{\text{mean}}$ (s) & TTFT (ms) & TPOT (ms) \\
\hline
0.6B, $k{=}4$, det   & \textbf{2.8548} & 0.3362 & 1.33 & 4.8881 & 132.57 & 41.23 \\
0.6B, $k{=}4$, stoch & 2.4798 & 0.2672 & 1.06 & 5.6196 & 134.54 & 45.51 \\
0.6B, $k{=}8$, det   & 1.9648 & 0.2378 & 1.86 & 7.1024 & 204.99 & 58.76 \\
0.6B, $k{=}8$, stoch & 1.6869 & 0.1916 & 1.50 & 8.2612 & 207.82 & 64.67 \\
0.6B, $k{=}16$, det  & 1.1091 & 0.1292 & 1.96 & 12.5817 & 349.58 & 94.24 \\
0.6B, $k{=}16$, stoch& 0.9719 & 0.1082 & 1.64 & 14.3390 & 354.82 & 106.03 \\
\hline
\end{tabular}
\end{table*}

\subsection{RQ1: Does speculative decoding accelerate inference?}

On RTX4090, all six Qwen2.5 settings and five of six Qwen3 settings exceed
baseline (Tables~\ref{tab:headline-4090-qwen25} and~\ref{tab:headline-4090}).
Both families peak at deterministic $k{=}4$ (Qwen2.5: $3.0674\times$; Qwen3:
$2.8548\times$). Qwen3 stochastic $k{=}16$ is the exception
($S{=}0.9719$), showing that long proposals can eliminate the gain.

\subsection{RQ2: How do $k$ and acceptance affect net latency?}

Larger $k$ increases $B_{\text{eff}}$ but reduces end-to-end speed, the
overhead inflection predicted by Eq.~\ref{eq:specdec-latency}. Mean
speedup falls from $2.9164$ to $1.8118$ for Qwen2.5 and from $2.6673$ to
$1.0405$ for Qwen3 between $k{=}4$ and $k{=}16$. Thus, additional drafting
and verification cost dominates the accepted-block benefit in this range.

\subsection{RQ3: Deterministic versus stochastic regime}

Qwen3 deterministic decoding is faster at every matched $k$, with its
advantage narrowing from $0.3750$ at $k{=}4$ to $0.1372$ at $k{=}16$.
Qwen2.5 follows the same pattern at $k{=}4$ and $8$, with a small reversal at
$k{=}16$. Deterministic decoding therefore provides the stronger low-$k$
runtime profile, although the preference is not universal.

\subsection{RQ4: Cross-hardware portability (RTX4090 vs RTX5090-laptop)}

\begin{table}[t]
\centering
\caption{Qwen3 fixed-policy results on RTX5090-laptop.
Ordering trends are preserved (best at low $k$, deterministic $>$ stochastic),
but absolute speedups are below 1 under the tested runtime stack.}
\label{tab:headline-5090}
\small
\begin{tabular}{l c c c}
\hline
\textbf{Config} & $S$ & $\alpha$ & $B_{\text{eff}}$ \\
\hline
0.6B, $k{=}4$, det   & 0.6743 & 0.3353 & 1.32 \\
0.6B, $k{=}4$, stoch & 0.5924 & 0.2845 & 1.06 \\
0.6B, $k{=}8$, det   & 0.4613 & 0.2538 & 1.85 \\
0.6B, $k{=}8$, stoch & 0.3995 & 0.2273 & 1.50 \\
0.6B, $k{=}16$, det  & 0.2588 & 0.1803 & 1.96 \\
0.6B, $k{=}16$, stoch& 0.2222 & 0.1543 & 1.64 \\
\hline
\end{tabular}
\end{table}

Cross-hardware validation (Table~\ref{tab:headline-5090}) reverses the practical outcome: all six
RTX5090-laptop settings are slower than their autoregressive baselines
($S{=}0.2222$--$0.6743$). Relative ordering remains consistent---low $k$ and
deterministic decoding perform best---but RTX4090 speedups do not transfer.
Deployment decisions must therefore be based on target-stack measurements,
not GPU generation or policy ranking alone.

Figure~\ref{fig:specdec-vs-baseline-metrics} consolidates the expanded
comparison across model families, hardware profiles, and repeat seeds.

\begin{figure*}[t]
\centering
\includegraphics[width=0.98\textwidth]{figures/specdec_configs_vs_baseline_metrics.pdf}
\caption{Expanded fixed-policy and stability comparison. (a) Qwen2.5 and
Qwen3 speedup on RTX4090; (b) Qwen3 speedup on RTX4090 and RTX5090-laptop;
(c) Qwen3 acceptance-rate portability; and (d) the change in mean latency
from repeat seed 123 to seed 999 at $k{=}4$ for both model families and
hardware profiles. Panels (a)--(c) use the primary seed-42 runs and
regime-matched autoregressive baselines. The dotted line in (a)--(b) marks
baseline parity ($S{=}1$).}
\label{fig:specdec-vs-baseline-metrics}
\end{figure*}

\subsection{Variability, token timing, and statistical tests}

For Qwen3 RTX4090, per-sample speedup (mean$\pm$std, $n{=}1000$) decreases
from $2.7561\pm0.5950$ to $1.3549\pm0.6533$ across deterministic
$k{=}4$--$16$, and from $2.4945\pm0.4565$ to $1.2173\pm0.6206$ under
stochastic decoding. Over the same range, mean TTFT rises from 132.57 to
349.58~ms (deterministic) and 134.54 to 354.82~ms (stochastic); TPOT rises
from 41.23 to 94.24~ms and 45.51 to 106.03~ms, respectively. Longer
proposals therefore degrade both initial and streaming responsiveness.

Seed stability was evaluated only at $k{=}4$ using the same manifests
($n{=}1000$) for seeds 123 and 999. Figure~\ref{fig:specdec-vs-baseline-metrics}(d)
summarizes the change in mean latency: deterministic runs differ by at most
1.5\%, while stochastic differences range from 0.4\% to 5.0\% across the two
families and devices. Regime ordering is preserved, and acceptance and token
timing vary less across seeds than across prompts. These repeats do not
establish seed stability for $k{=}8$ or $16$.

One-sided paired $t$-tests ($H_0{:}\ S\leq1$), paired by sample and corrected
across six Qwen3 RTX4090 configurations using Holm--Bonferroni, confirm
speedup for deterministic $k{=}4$, $8$, and $16$ ($d_S{=}1.18$, $1.10$, and
$0.29$; all $p_S<10^{-18}$) and stochastic $k{=}4$ and $8$
($d_S{=}1.24$ and $1.13$; both $p_S<10^{-179}$). Stochastic
$k{=}16$ is not significant ($d_S{=}-0.08$, $p_S{=}0.996$),
consistent with its below-baseline mean. Tests use the seed-42 full grid;
repeat seeds cover $k{=}4$ only.

CNN/DailyMail quality also changes under deterministic speculative decoding
($n{=}200$). Relative to the autoregressive baseline
($Q_{\mathrm{sum}}{=}19.58$, $Q_{\mathrm{bleu}}{=}3.20$,
$Q_{\mathrm{bert}}{=}0.793$), $k{=}4$ yields 11.50, 1.39, and 0.710;
$k{=}16$ partially recovers to 14.92, 2.04, and 0.741. All three metrics
therefore identify a speed--quality trade-off rather than free acceleration.

\subsection{Hypothesis resolution}

H1 is supported for all Qwen2.5 and five of six Qwen3 RTX4090 settings, but
rejected for the RTX5090-laptop matrix. H2 is rejected: increasing $k$ from
4 to 8 reduces speedup. H3 is supported because $B_{\text{eff}}$ grows while
latency worsens. H4 is mostly supported: deterministic decoding leads at all
Qwen3 points and at low-to-mid $k$ for Qwen2.5, but Qwen2.5 $k{=}16$ shows a
small reversal. These outcomes bound the claims to the evaluated model,
hardware, and runtime combinations.
